\documentclass[11pt]{article}
\usepackage{fontspec}
\setmainfont{Times New Roman}
\setsansfont{Arial}
\setmonofont{Consolas}
\usepackage{amsmath,amssymb}
\usepackage{geometry}
\usepackage{microtype}
\geometry{a4paper,margin=3cm}
\setlength{\parindent}{1.25em}
\setlength{\parskip}{0pt}
\emergencystretch=30em
\allowdisplaybreaks
\raggedbottom
\newcommand{\ideal}[1]{\mathfrak{#1}}

\begin{document}

% Unit: Paper 19, Section 3. Deterministic Interslavic Cyrillic reader generated from the Latin source and manually checked.

\setcounter{footnote}{12}
\subsection*{\S~3. Равност числа компонентов при двох различних разложеньах на неразложиме идеалы.}

За доказ равности числа треба најпрво изразити разложимост односно неразложимост идеала чрез својства јего комплемента по следујучем теорему.

\textbf{Теорем III.}\footnote{Теорем III одговарја преходу од модулов к фактор-групам в работах Schmeidlera и Noether--Schmeidlera (ср. ввод). Идеалу \(\ideal C\) одговарја фактор-група, а \(\ideal N_1\) и \(\ideal N_2\) подгрупы, в кторе фактор-група разлагаје се.} \emph{Најкратше представјенье \(\ideal M=[\ideal U,\ideal C]\) нај буде редуковано в односу к \(\ideal C\). Тогда нужно и достаточно условје за то, же \(\ideal C\) јест разложимы, јест сушчствованье двох идеалов \(\ideal N_1\) и \(\ideal N_2\), кторе сут властне делители од \(\ideal M\), тако же}
\begin{equation}
 \ideal N_1\ideal N_2\equiv0(\ideal M),\qquad
 \ideal N_i\equiv0(\ideal U),\qquad
 [\ideal N_1,\ideal N_2]=\ideal M.
\tag{2}
\end{equation}
\emph{Из того још следује: ако условја (2) сут изполньене, а \(\ideal C\) јест неразложимы, тогда најменеј једно \(\ideal N_i\) не јест властны делитељ од \(\ideal M\); \(\ideal N_i=\ideal M\).}

Нај буде \(\ideal C=[\ideal C_1,\ideal C_2]\), гдје \(\ideal C_1,\ideal C_2\) сут властне делители од \(\ideal C\). Тогда выходи
\[
 \ideal M=[\ideal U,\ideal C]=[\ideal U,\ideal C_1,\ideal C_2]
       =[[\ideal U,\ideal C_1],[\ideal U,\ideal C_2]].
\]
Ту идеалы \([\ideal U,\ideal C_i]\) сут властне делители од \(\ideal M\), бо иначе \([\ideal U,\ideal C]\) не было бы редуковано в односу к \(\ideal C\). Понеже такоже делимост чрез \(\ideal U\) јест изполньена, условје (2) јест доказано како нужно. (Представјенье (2) не јест редуковано, бо једно \([\ideal U,\ideal C_i]\) може быти заменьено чрез \(\ideal C_i\).)

Нај тепер обратно (2) буде изполньено. Творимо идеалы
\[
 \ideal C_1=(\ideal C,\ideal N_1),\qquad
 \ideal C_2=(\ideal C,\ideal N_2),\qquad
 \ideal C^*=[\ideal C_1,\ideal C_2].
\]
Тогда \(\ideal C\) јест делимы и чрез \(\ideal C_1\), и чрез \(\ideal C_2\), зато такоже чрез најменше обче многократно \(\ideal C^*\). Да бы се показала делимост од \(\ideal C^*\) чрез \(\ideal C\), нај
\[
 f\equiv0(\ideal C^*);\quad f\equiv0(\ideal C_1);\quad f\equiv0(\ideal C_2),
 \quad\text{или такоже}\quad f=c+n_1=t+n_2,
\]
гдје \(c,t,n_1,n_2\) сут величины из \(\ideal C,\ideal C,\ideal N_1,\ideal N_2\), значи особливо \(n_1\) јест делимы чрез \(\ideal N_1\). Зато разлика
\[
 g=c-t=n_2-n_1
\]
јест делимы и чрез \(\ideal C\), и чрез \(\ideal U\), зато чрез \(\ideal M\). По \(n_1=n_2+m\) јест далье \(n_1\) (такоже \(n_2\)) делимы и чрез \(\ideal N_1\), и чрез \(\ideal N_2\), зато чрез \(\ideal M\). Зато става
\[
 f=c+m,\qquad f\equiv0(\ideal C),\qquad \ideal C^*=\ideal C.
\]
\(\ideal C_1\) и \(\ideal C_2\) сут при том властне делители од \(\ideal C\); бо при \(\ideal C_1=(\ideal C,\ideal N_1)=\ideal C\) был бы \(\ideal N_1\) делимы чрез \(\ideal C\), зато из делимости чрез \(\ideal U\) был бы ровны \(\ideal M\), против предпоставце. Зато \(\ideal C=[\ideal C_1,\ideal C_2]\) јест познаны како разложимы; теорем III јест доказан.

Треба приметити, же скоро то само размысльенье показује још следујуче:

\textbf{Лема II.} \emph{Ако в најкратшем представјеньу \(\ideal M=[\ideal U,\ideal C]\) идеал \(\ideal C\) може быти заменьен властным делительем, тогда \(\ideal C\) јест разложимы.}

Нај зато
\[
 \ideal M=[\ideal U,\ideal C]=[\ideal U,\ideal C_1],
\]
и нај буде положено
\[
 \ideal C^*=[\ideal C_1,(\ideal U,\ideal C)].
\]
Тогда \(\ideal C\) зново јест делимы чрез \(\ideal C^*\). Из \(f\equiv0(\ideal C^*)\) далье следује
\[
 f=c_1=a+c.
\]
Разлика \(a=c_1-c\) јест зато делимы и чрез \(\ideal U\), и чрез \(\ideal C_1\), следовательно чрез \(\ideal M\); зато става \(c_1=c+m\), \(f\equiv0(\ideal C)\), \(\ideal C^*=\ideal C\). Понеже и \(\ideal C_1\), и \((\ideal U,\ideal C)\) сут по предпоставце властне делители од \(\ideal C\), тым \(\ideal C=\ideal C^*\) јест познаны како разложимы.

Неразложимы \(\ideal C\) зато не може быти заменьен властным делительем.

Нај тепер будут дане две различне најкратше представјеньа од \(\ideal M\) како најменшего обчего многократного конечно многых неразложимых идеалов:
\[
 \ideal M=[\ideal B_1\cdots\ideal B_k]=[\ideal D_1\cdots\ideal D_l].
\]
Ты представјеньа сут по приметце к леми II такисто редуковане. Тогда најпрво доказујемо

\textbf{Лема III.} \emph{К всакому комплементу \(\ideal U_i=[\ideal B_1\cdots\ideal B_{i-1}\ideal B_{i+1}\cdots\ideal B_k]\) сушчствује идеал \(\ideal D_j\), тако же \(\ideal M=[\ideal U_i,\ideal D_j]\).}

Именно, ако положимо \(\ideal M=[\ideal B_i,\ideal C_1]\), \(\ideal C_1=[\ideal D_1,\ideal C_2]\) итд., тогда добива се
\[
 \ideal M=[\ideal U_i,\ideal B_i]=[\ideal U_i,[\ideal D_1,\ideal C_2]]=[\ideal B_i,[\ideal U_i,\ideal D_1],\ideal C_2].
\]
Ту за \(\ideal N_1=[\ideal U_i,\ideal D_1]\), \(\ideal N_2=[\ideal U_i,\ideal C_2]\) сут изполньене условја (2) теорема III, понеже \(\ideal M=[\ideal U_i,\ideal B_i]\) јест редуковано в односу к \(\ideal B_i\), а представјенье јест најкратше. Понеже \(\ideal B_i\) был предпоставјен како неразложимы, нужно једно \(\ideal N_i\) мора стати ровно \(\ideal M\).

Ако \([\ideal U_i,\ideal D_1]=\ideal M\), лема јест доказана; ако \([\ideal U_i,\ideal C_2]=\ideal M\), одповедајуче добива се \(\ideal M=[[\ideal U_i,\ideal D_1],[\ideal U_i,\ideal C_2]]\), гдје по том самом закльучку зново једна компонента мора стати ровна \(\ideal M\). Тако продолжајучи, долази или \(\ideal M=[\ideal U_i,\ideal D_j]\), гдје \(j<l\); или става \(\ideal M=[\ideal U_i,\ideal D_1,\ldots,\ideal D_{l-1}]\); по \(\ideal D_1\cdots\ideal D_{l-1}=\ideal D_l\) тым лема јест доказана.

Из того тепер следује

\textbf{Теорем IV.} \emph{При двох различних најкратших представјеньах идеала како најменшего обчего многократного из неразложимых, число компонентов јест исто.}

Из лемы именно следује за \(i=1\):
\[
 \ideal M=[\ideal B_1\cdots\ideal B_k]=[\ideal U_1,\ideal D_{j_1}]
 =[\ideal D_{j_1},\ideal B_2,\ldots,\ideal B_k].
\]
Ако тепер разгледајемо две разложеньа
\[
 \ideal M=[\ideal D_{j_1},\ideal B_2,\ldots,\ideal B_k]
          =[\ideal D_1,\ideal D_2,\ldots,\ideal D_l],
\]
и повторимо горньи закльучак в односу к комплементу \(\ideal U_2=[\ideal D_{j_1},\ideal B_3,\ldots,\ideal B_k]\) од \(\ideal B_2\), тогда добива се
\[
 \ideal M=[\ideal U_2,\ideal B_2]=[\ideal U_2,\ideal D_{j_2}]
          =[\ideal D_{j_1},\ideal D_{j_2},\ideal B_3,\ldots,\ideal B_k],
\]
и продолженьем процедуры
\[
 \ideal M=[\ideal D_{j_1},\ideal D_{j_2},\ldots,\ideal D_{j_k}].
\]
Понеже по предпоставце представјенье \(\ideal M=[\ideal D_1\cdots\ideal D_l]\) јест најкратше, значи ниједно \(\ideal D\) не може быти опушчено, различне медју \(\ideal D_{j_i}\) морајут исцрпати все \(\ideal D\); зато \(k\ge l\). Ако в леми и следујучих закльучках всуды заменимо \(\ideal B\) с \(\ideal D\), одповедајуче добива се \(l\ge k\), и зато \(k=l\), чиме равност числа јест доказана. Из того још следује, же идеалы \(\ideal D_{j_i}\) сут вси меджу собоју различне, бо иначе в најкратшем представјеньу чрез \(\ideal D_{j_i}\) појавило бы се менье неж \(k\) компонентов; можно зато выбрати означенье тако, же \(\ideal D_{j_i}=\ideal D_i\). Из тој самој причины сут такоже вса меджупредставјеньа \(\ideal M=[\ideal D_{j_1}\cdots\ideal D_{j_r},\ideal B_{r+1},\ldots,\ideal B_k]\) најкратше и по приметце к леми II зато редуковане.

Равност числа води к обрату лемы I чрез

\textbf{Лема IV.} \emph{Ако се в редукованом представјеньу компоненты схвате в групы и створи се њихово најменше обче многократно, наставајуче представјенье буде редуковано. Другыми словами: из редукованого представјеньа}
\[
 \ideal M=[\ideal C_1\cdots\ideal C_k]
\]
\emph{следује, же такоже \(\ideal M=[\ideal N_1,\ideal N_2]\) јест редуковано, ако јест положено \(\ideal N_1=[\ideal C_1,\,\ldots,\ideal C_h]\), \(\ideal N_2=[\ideal C_{h+1},\ldots,\ideal C_k]\).}

Најпрво треба приметити, же \(\ideal N_1\) не може делити свој комплемент \(\ideal N_2\), бо то не важи за ниједного из јего делительев \(\ideal C_i\); представјенье јест зато најкратше. Да бы се показало, же \(\ideal N_1\) не може быти заменьено нијакым властным делительем, разлагајемо \(\ideal C\) на јих неразложиме идеалы \(\ideal B\),\footnote{Под тым треба всегда разумети најкратше, значи редуковано представјенье чрез \(\ideal B\).} тако же наставајут (по леми I) редуковане представјеньа
\[
 \ideal N_1=[\ideal B_1\cdots\ideal B_\lambda],
 \qquad
 \ideal N_2=[\ideal B_{\lambda+1}\cdots\ideal B_\kappa].
\]
Нај тепер \(\ideal M=[\ideal N_1,\ideal N_2]\) буде редуковано в односу к \(\ideal N_2\), а \(\ideal N_1^*\) властны делитель од \(\ideal N_1\). По леми II става
\[
 \ideal N_1=[\ideal N_1^*,(\ideal N_1,\ideal N_2)],
\]
и то представјенье јест нужно редуковано в односу к \(\ideal N_1^*\), бо иначе \(\ideal N_1\) могло бы се и в \(\ideal M\) заменити властным делительем. Заменимо тепер по потребе такоже \((\ideal N_1,\ideal N_2)\) властным делительем, тако же настане редуковано представјенье за \(\ideal N_1\). Ако тепер разложимо обе компоненты од \(\ideal N_1\) на неразложиме идеалы, число \(\lambda\) неразложимых идеалов од \(\ideal N_1\) склада се адитивно из чисел компонентов; число неразложимых идеалов одговарјајучих властному делительу \(\ideal N_1^*\) јест зато нужно менье неж \(\lambda\). Тогда але и разложенье од \(\ideal M=[\ideal N_1^*,\ideal N_2]\) на неразложиме идеалы води к менье неж \(\kappa\) идеалам, в противречности с равностју числа. Како специалны случај \(\varrho=2\) још следује, же представјенье \(\ideal M=[\ideal N_1,\ideal N_2]\) јест такоже редуковано в односу к комплементу \(\ideal N_2\).

\clearpage

\end{document}

